\documentclass[11pt]{article}

\usepackage[margin=1in]{geometry}
\usepackage{amsmath,amssymb,amsthm}
\usepackage{mathtools}

\newtheorem{theorem}{Theorem}
\newtheorem{proposition}{Proposition}
\newtheorem{lemma}{Lemma}
\theoremstyle{remark}
\newtheorem{remark}{Remark}

\DeclareMathOperator{\msp}{mp}
\newcommand{\content}{\operatorname{content}}
\newcommand{\stil}{\widetilde{s}}
\newcommand{\ones}[1]{1^{#1}}

% \author{}  % add author line if desired
\title{The $\stil_{\varnothing}$-coefficient of $s_{(n-1,1)}$}
\date{}

\begin{document}
\maketitle

\noindent
Throughout, $p(k)$ denotes the number of integer partitions of $k$, with the
convention $p(0)=1$, and for a finite multiset $M$ we write $\msp(M)$ for the
number of multiset partitions of $M$ (unordered collections of nonempty
sub-multisets whose union is $M$). We work in the ring of symmetric functions
with $\{\stil_\mu\}$ the Orellana--Zabrocki irreducible character basis
\cite{OZ}. For a partition $\lambda=(\lambda_1,\lambda_2,\dots)$ the
\emph{content} of $h_\lambda$ is the multiset
$\content(\lambda)=\{\!\{\ones{\lambda_1},2^{\lambda_2},\dots\}\!\}$, i.e.\ it
contains $\lambda_i$ copies of the label $i$. For partitions $\lambda,\mu$ we let
$N_{\lambda,\mu}$ be the coefficient of $\stil_\mu$ in the expansion of the Schur
function $s_\lambda$.

The goal is a closed form for $N_{(n-1,1),\varnothing}$ together with a proof
that it is nonnegative.

\section{Two recalled facts}

\begin{lemma}\label{lem:empty}
For every partition $\lambda$,
\[
  \bigl[\stil_{\varnothing}\bigr]\,h_\lambda \;=\; \msp\bigl(\content(\lambda)\bigr).
\]
\end{lemma}

\begin{proof}
By \cite[Prop.~2]{OZ}, for any fixed $m\ge 0$ the coefficient of $\stil_\mu$ in
$h_\lambda$ counts the column-strict tableaux with $m$ blank cells in the first
row, all remaining cells filled with nonempty multisets of total content
$\content(\lambda)$, and body shape (the shape below the first row) equal to
$\mu$. For $\mu=\varnothing$ there are no cells below the first row, so such a
tableau is a single row: $m$ blanks followed by nonempty multiset cells that are
weakly increasing in the lexicographic order on multisets. Discarding the blanks,
the filled cells are precisely a multiset partition of $\content(\lambda)$
written in its unique weakly increasing order. This assignment is a bijection
between these tableaux and the multiset partitions of $\content(\lambda)$, whence
the count is $\msp(\content(\lambda))$.
\end{proof}

\begin{lemma}\label{lem:jt}
For $n\ge 2$,
\[
  s_{(n-1,1)} \;=\; h_{(n-1,1)} - h_{(n)}.
\]
\end{lemma}

\begin{proof}
The two-row Jacobi--Trudi identity gives
\[
  s_{(n-1,1)}
  = \det\begin{pmatrix} h_{n-1} & h_{n}\\ h_{0} & h_{1}\end{pmatrix}
  = h_{n-1}h_{1} - h_{n}h_{0}
  = h_{(n-1,1)} - h_{(n)},
\]
using $h_0=1$.
\end{proof}

\section{The closed form}

\begin{proposition}\label{prop:count}
For $n\ge 1$,
\[
  \msp\bigl(\{\!\{\ones{n-1},2\}\!\}\bigr) \;=\; \sum_{k=0}^{n-1} p(k).
\]
\end{proposition}

\begin{proof}
Fix a multiset partition of $\{\!\{\ones{n-1},2\}\!\}$. Exactly one of its blocks
contains the single label $2$; say that block contains $j$ ones as well, where
$0\le j\le n-1$. The remaining $n-1-j$ ones are partitioned among the other
blocks, and since these ones are indistinguishable this is an integer partition
of $n-1-j$, contributing $p(n-1-j)$ possibilities. Conversely any choice of $j$
and of a partition of $n-1-j$ determines the multiset partition. Summing over
$j$,
\[
  \msp\bigl(\{\!\{\ones{n-1},2\}\!\}\bigr)
  = \sum_{j=0}^{n-1} p(n-1-j)
  = \sum_{k=0}^{n-1} p(k). \qedhere
\]
\end{proof}

\begin{theorem}\label{thm:formula}
For $n\ge 2$,
\[
  N_{(n-1,1),\varnothing} \;=\; \sum_{k=0}^{n-1} p(k)\;-\;p(n).
\]
\end{theorem}

\begin{proof}
Apply $[\stil_\varnothing]$ to Lemma~\ref{lem:jt} and use Lemma~\ref{lem:empty}:
\[
  N_{(n-1,1),\varnothing}
  = \bigl[\stil_\varnothing\bigr]h_{(n-1,1)} - \bigl[\stil_\varnothing\bigr]h_{(n)}
  = \msp\bigl(\{\!\{\ones{n-1},2\}\!\}\bigr) - \msp\bigl(\{\!\{\ones{n}\}\!\}\bigr).
\]
Now $\msp(\{\!\{\ones{n}\}\!\})=p(n)$, since a multiset partition of $n$
indistinguishable ones is exactly an integer partition of $n$; and
$\msp(\{\!\{\ones{n-1},2\}\!\})=\sum_{k=0}^{n-1}p(k)$ by
Proposition~\ref{prop:count}. Substituting gives the claim.
\end{proof}

\section{Nonnegativity via an injection}

Let $\mathcal{P}(n)$ be the set of integer partitions of $n$, identified with the
multiset partitions of $\{\!\{\ones{n}\}\!\}$ (a partition
$\lambda=(\lambda_1\ge\cdots\ge\lambda_\ell)$ corresponds to the block collection
$\{\!\{\{\!\{\ones{\lambda_1}\}\!\},\dots,\{\!\{\ones{\lambda_\ell}\}\!\}\}\!\}$),
and let $\mathcal{Q}(n)$ be the set of multiset partitions of
$\{\!\{\ones{n-1},2\}\!\}$. Thus $|\mathcal{P}(n)|=p(n)$ and
$|\mathcal{Q}(n)|=\sum_{k=0}^{n-1}p(k)$ by Proposition~\ref{prop:count}.

\begin{theorem}\label{thm:inj}
There is an injection $\iota:\mathcal{P}(n)\hookrightarrow\mathcal{Q}(n)$.
Consequently $N_{(n-1,1),\varnothing}=|\mathcal{Q}(n)|-|\mathcal{P}(n)|\ge 0$.
\end{theorem}

\begin{proof}
\emph{Definition of $\iota$.} Given $\lambda\in\mathcal{P}(n)$, view it as the
above block collection and replace one block of maximum size, namely a block
$\{\!\{\ones{\lambda_1}\}\!\}$, by $\{\!\{\ones{\lambda_1-1},2\}\!\}$. That is,
in one largest block, turn a single $1$ into the label $2$. The union of the
blocks becomes $\{\!\{\ones{n-1},2\}\!\}$, so $\iota(\lambda)\in\mathcal{Q}(n)$.

\emph{Well-defined.} If $\lambda$ has several parts equal to $\lambda_1$, the
blocks $\{\!\{\ones{\lambda_1}\}\!\}$ are equal as multisets and hence
interchangeable within the (unordered) multiset partition; modifying any one of
them yields the same element of $\mathcal{Q}(n)$. So $\iota(\lambda)$ does not
depend on the choice.

\emph{Injective.} In $\iota(\lambda)$ exactly one block contains the label $2$;
write it as $\{\!\{\ones{s},2\}\!\}$. By construction $s=\lambda_1-1$ and every
other block is one of the $\{\!\{\ones{\lambda_i}\}\!\}$ with $2\le i\le\ell$.
Replacing the block $\{\!\{\ones{s},2\}\!\}$ by $\{\!\{\ones{s+1}\}\!\}$ recovers
the block collection of $\lambda$, hence $\lambda$ itself. Thus $\lambda$ is
determined by $\iota(\lambda)$, and $\iota$ is injective.

The inequality $N_{(n-1,1),\varnothing}\ge 0$ then follows from
Theorem~\ref{thm:formula}, since an injection forces
$|\mathcal{P}(n)|\le|\mathcal{Q}(n)|$.
\end{proof}

\begin{remark}[The coefficient counts the cokernel]
Write each element of $\mathcal{Q}(n)$ by isolating its unique block
$\{\!\{\ones{s},2\}\!\}$ containing the label $2$. The image of $\iota$ is exactly
those elements for which $s+1\ge\mu$ for every all-ones block
$\{\!\{\ones{\mu}\}\!\}$ present, i.e.\ those in which the $2$-block, after its
$2$ is turned back into a $1$, is a block of maximum size. Therefore
\[
  N_{(n-1,1),\varnothing}
  = \#\Bigl\{\,\pi\in\mathcal{Q}(n)\ :\ \text{some all-ones block of }\pi
      \text{ has strictly more than } s+1 \text{ ones}\,\Bigr\},
\]
where $\{\!\{\ones{s},2\}\!\}$ is the $2$-block of $\pi$. This realizes the
coefficient as the cardinality of the complement of $\operatorname{im}\iota$.
\end{remark}

\begin{remark}[Small values and consistency]
For $n=2$, $N_{(1,1),\varnothing}=p(0)+p(1)-p(2)=1+1-2=0$, and indeed $\iota$ is a
bijection ($\mathcal{Q}(2)=\{\{\!\{1,2\}\!\},\,\{\!\{1\}\!\}\{\!\{2\}\!\}\}$, both
in the image). For $n=3$, $N_{(2,1),\varnothing}=(1+1+2)-3=1$, matching the
coefficient of $\stil_\varnothing$ in the expansion
$s_{(2,1)}=\stil_\varnothing+3\stil_{(1)}+2\stil_{(2)}+2\stil_{(1,1)}+\stil_{(2,1)}$;
the single element outside $\operatorname{im}\iota$ is
$\{\!\{1,1\}\!\}\{\!\{2\}\!\}$. Further values: $N=2,5,8$ for $n=4,5,6$.
\end{remark}

\begin{remark}[Relation to the injection program]
This is the $b=1$ instance of the general target
$M(a+1,b-1)\hookrightarrow M(a,b)$, where $M(a,b)=\msp(\{\!\{1^{a}2^{b}\}\!\})$: for
$b=1$ and $a=n-1$ it reads $M(n,0)\hookrightarrow M(n-1,1)$, and the map $\iota$
above is an explicit such injection.
\end{remark}

\begin{thebibliography}{9}
\bibitem{OZ}
R.~Orellana and M.~Zabrocki,
\emph{Symmetric group characters as symmetric functions},
Adv.\ Math.\ \textbf{390} (2021), Paper No.\ 107943.
\end{thebibliography}

\end{document}